\section{Conclusion and Future Directions}

\subsection{Summary of Contributions}
{\sloppy
This research successfully develops a Coarse\-Grained Reconfigurable Architecture (CGRA) tailored for sparse SNN inference and convolutional neural network acceleration. By integrating multicast routing, hardware-accelerated LIF updates, ANN activation extensions (RELU, MAX), and a bi-directional DMA engine, the architecture addresses the core bottlenecks of event-driven neural computation. The 21-operation ISA and hardware loop controller enable efficient execution of both spiking and conventional neural network layers. A complete software stack---including a UIO+CMA Linux driver, Im2Col tiler library, and License Plate Recognition demo application---demonstrates end-to-end deployment readiness on the Xilinx Zynq-7000 XC7Z020 (Haoyue Zynq-7020 board, PetaLinux 2022.2). The 100\% pass rate across 8,915 regression tests (25 suites) and timing closure at 50 MHz on the XC7Z020 (post-synthesis headroom of 213.72 MHz on UltraScale+ reference) prove the feasibility and robustness of the design. The full hardware/software implementation is available as open-source documentation \cite{cgra_project_code}.
}

\subsection{Future Directions}
While the current architecture provides a solid baseline for SNN inference, several areas for future development are identified:
\begin{enumerate}
    \item \textbf{On-chip Learning:} Integration of unsupervised learning rules such as Spike-Timing-Dependent Plasticity (STDP) directly into the PE datapath.
    \item \textbf{Scalability:} Extending the 4×4 mesh to larger configurations (e.g., 8×8 or 16×16) and evaluating hierarchical interconnect strategies.
    \item \textbf{Configurable Precision:} Supporting mixed-precision levels (e.g., 4-bit, 8-bit, 16-bit) to allow trading off accuracy for throughput and energy efficiency.
    \item \textbf{Compiler Toolchain:} Development of an automated mapping tool that translates SNN models from high-level frameworks (PyTorch, SNN-lib) to CGRA configuration streams.
\end{enumerate}
